\section{Preparing an ordered half-deck}\label{sec:preparation}

Throughout this section put $m=n/2$. Our first task is to randomize any specified ordered half-deck from a deterministic start in $400$ sweeps, uniformly in its starting positions. The conditional-bijection construction of Lemma~\ref{lem:slot-kernel} supplies its one-card transition matrices. We first establish the complementary one-sweep estimates that will make these matrices contract over pairs of sweeps.

Consider one sweep with blocked positions \(E\) at its start, \(r=|E|\le m\). For the random blocked paths let \(E'\) be the end positions and \(Q\) the conditional single-free-card transition matrix (rows in \(E^c\), columns in \((E')^c\)). It is doubly stochastic between these two sets, since the conditional moves are random bijections. Vectors of probabilities or their differences will be row vectors, with the ordinary Euclidean norm. In particular \(\|uQ\|\le\|u\|\), by convexity of the square applied in each column and then the row sums. Put \(A_E=\mathbb E[QQ^\top]\), the expectation over these paths from the fixed blocked start. Its entry at \(x,y\in E^c\) is the probability two copies starting at \(x,y\), run independently conditional on the common blocked paths, finish at the same position (averaging also over the paths). It is symmetric and a contraction, as is \(QQ^\top\).

Let \(h=\lfloor d/2\rfloor\), and call the blocks specified by bits \(>h\) the coarse blocks (size \(2^h\)). For a free-slot vector \(u\) write \(u_{\rm c}\) for its orthogonal projection obtained by averaging over free slots within each such block (ignoring empty ones). Call a blocked set good if its fraction in every coarse block is at most \(3/4\). Set \(q=7/8\) and
\[
 \epsilon_n=n\left(\frac{3/2}{2^{3/4}}\right)^{2^h}.
\]

We use a two-sweep mechanism. From a balanced blocked set, one sweep contracts fluctuations within large subcubes. It may leave averages that differ between subcubes, but the same sweep makes the expected energy of those output averages small. Combining the two estimates in the next sweep will contract the whole vector. The next lemma gives the two estimates and the probability of an unbalanced output.

\begin{lemma}[Coarse-block estimates]\label{lem:coarse-block-estimates}
For the one-sweep matrix $Q$ defined above and every free-slot row vector $u$, a good starting blocked set $E$ satisfies
\begin{equation}
 \mathbb E\|uQ\|^2=u A_E u^\top
 \le \|u_{\rm c}\|^2+q^{d-h}\|u\|^2. \label{special:burn-active:eq:2}
\end{equation}
For every starting blocked set $E$ with $|E|\le n/2$, the ending set satisfies
\begin{equation}
 \mathbb P(E'\text{ not good})\le
 \epsilon_n. \label{special:burn-active:eq:3}
\end{equation}
If $u$ has sum zero, then its coarse part on the ending free slots satisfies
\begin{equation}
 \mathbb E\|(uQ)_{\rm c}\|^2\le (4\cdot 2^{-h}+\epsilon_n)\|u\|^2. \label{special:burn-active:eq:4}
\end{equation}
\end{lemma}
\begin{proof}
Write \(p_k(x)\) for the initial blocked fraction in the block of size \(2^{k-1}\) having bits \(>k\) equal to those of \(x\) and bit \(k\) opposite to that of \(x\), and set \(q_k(x)=(1+p_k(x))/2\). Let \(j\) be the largest coordinate where \(x,y\) differ, or \(0\) if they coincide. We claim
\begin{equation}
 A_E(x,y)=2^{-j}\prod_{k=j+1}^d q_k(x). \label{special:burn-active:eq:1}
\end{equation}
For the two-copy experiment averaged over paths, we can generate the blocked motion with independent switch coins, use these moves for the copies where an edge is incident to a blocked card, and use two fresh independent arrays of switch coins (one per copy) on free-free edges. Whenever the copies are at different positions before a stage, their two new bits at that stage are independent uniform bits given the history up to then: for different edges the coins used are independent, and if the edge is the same, both its slots must be free and the copies use independent coins. Up to and including stage \(j\) the pre-switch positions are different, due to bit \(j\) if \(j>0\). Thus the probability of matching the first \(j\) new bits is \(2^{-j}\), after which they are at the same position on that event. A mismatch in any updated bit persists to the end of the sweep.

For \(k>j\), given they matched the new bits through stage \(k-1\) (so are now together), they stay together for sure if the neighbor slot is blocked and with probability \(1/2\) otherwise. Here the probability the neighbor is blocked, given this matching so far and their current common prefix, is \(p_k(x)\). To see this, before stage \(k\) both copies have evolved in the block with bits \(\ge k\) equal to those of \(x\). Their motion and the event of matching so far only use coin arrays in that block. In the neighboring block (bit \(k\) opposite), independently, the blocked configuration has evolved for \(k-1\) stages within the block; each initial blocked card there has uniform marginal position in it, since its successive new bits are conditionally fair. Thus the blocked probability at the slot matching the copies' prefix is the stated fraction. This independence uses a fixed initial \(E\) and disjoint switches of the blocks prior to stage \(k\). Multiplying the stay-together probabilities proves \eqref{special:burn-active:eq:1}, including for \(j=0\).

To prove \eqref{special:burn-active:eq:2}, note that \eqref{special:burn-active:eq:1} is constant between any two fixed distinct coarse blocks. Within the block of \(x\), its row sum equals \(\prod_{k=h+1}^d q_k(x)\): for \(1\le j\le h\) there are \(2^{j-1}(1-p_j(x))\) free \(y\)'s with that last difference, contributing \((1-q_j(x))\prod_{k>j}q_k(x)\), which telescope together with the diagonal term. This row sum is constant for \(x\) in the block, and is at most \(q^{d-h}\) when \(E\) is good since the blocks defining \(p_k\) for \(k>h\) are unions of coarse blocks. The coarse-constant subspace is invariant under \(A_E\) by the constants and row sums just noted; on it we use contraction. On its orthogonal complement the between-block entries contribute zero, and the symmetric within-block matrices have norm at most their maximum row sum (nonnegative entries). By symmetry there is no cross term between the two subspaces. This proves \eqref{special:burn-active:eq:2}; without goodness we still have contraction.

We also bound the expected coarse part at the \emph{output}, with the same choice of coarse coordinates for use in the next sweep. Let \(H(x,y)\) be the probability in the two-copy experiment of equal bits \(>h\) at the output, and \(H_{\rm pre}(x,y)\) the probability of equal bits \(1,\ldots,h\). Then, with the all-ones matrix denoted by \({\bf 1}{\bf 1}^\top\),
\[
 H=2^{-(d-h)}({\bf 1}{\bf 1}^\top-H_{\rm pre})+A_E .
\]
If the first \(h\) bits mismatch, the copies remain at distinct positions, so the remaining new bits match with probability \(2^{-(d-h)}\) by the fresh-bit observation for distinct positions. If the first \(h\) match, also matching the rest means an end collision. Moreover \(H_{\rm pre}\) is positive semidefinite: conditional on the blocked paths the equality probabilities form a Gram matrix of the transition probabilities aggregated by output prefix, and we average this matrix. Consequently for \(u\) of sum zero, the expected sum over coarse output blocks of the squared sums of \(uQ\) in them is \(u H u^\top\le u A_E u^\top\le\|u\|^2\).

We prove directly the needed negative-dependence inequality. It is a restricted form of the preservation of negative dependence under partial symmetrization studied by Borcea, Br\"and\'en and Liggett~\cite[Theorems~4.9 and~4.20]{bbl2009}. The blocked-slot indicators throughout the switches satisfy: for any set of slots the probability all are blocked is bounded by the product of their marginal probabilities. It holds initially (deterministic); a single fair swap preserves it. For a test set containing one of the two slots, average the two old bounds, and for one containing both, use that the product of their marginals can only increase when both marginals are replaced by their average. Disjoint simultaneous swaps can be processed one at a time. After a sweep all marginals are \(r/n\) by single-card uniformity. For the number \(B\) blocked in a given coarse block, expanding the product of \(1+\)indicator over it thus gives \(\mathbb E 2^B\le(1+r/n)^{2^h}\); Markov's inequality and the union bound give \eqref{special:burn-active:eq:3}.

On good output the squared block sums are divided by block free-slot counts at least \(2^h/4\) to form the squared norm (use the unconditional numerator bound), and on bad output the norm is at most \(\|u\|\) by contraction. Together with \eqref{special:burn-active:eq:3}, this proves \eqref{special:burn-active:eq:4}.

\end{proof}

\begin{proposition}[Uniform half-deck preparation]\label{prop:half-preparation}
For all sufficiently large $d$, every ordered list of $k\le n/2$ distinct labels, from every deterministic starting layout, has endpoint law within $n^{-3}$ in total variation of uniform ordered distinct positions after $400$ sweeps.
\end{proposition}
\begin{proof}
Fix at most $m$ blocked cards and track one specified free card over sweeps, starting deterministically, and let \(u_s\) be its law conditional on the blocked paths up through sweep \(s\), minus uniform on the then free slots. Given those paths, the next blocked paths are fresh from the then layout, and on revealing them \(u_{s+1}=u_s Q\) for their conditional matrix; future blocked paths supply no information on the unused free coins of the past, and the uniform vectors are carried to uniform vectors. In particular there is pathwise norm contraction. Write \(a_s=\mathbb E\|u_s\|^2\), with \(a_0\le1\). For \(s\ge1\), the mean squared coarse part of \(u_s\) is bounded by \((4\cdot 2^{-h}+\epsilon_n)a_{s-1}\) by \eqref{special:burn-active:eq:4} given the paths through the previous sweep. Also \(\mathbb E[{\bf 1}_{\{\text{bad after }s\}}\|u_s\|^2]\le\epsilon_n a_{s-1}\) by pathwise contraction and \eqref{special:burn-active:eq:3} given the previous paths. Applying \eqref{special:burn-active:eq:2} for the next sweep on good starts and contraction otherwise, we get
\[
 a_{s+1}\le q^{d-h}a_s+(4\cdot 2^{-h}+2\epsilon_n)a_{s-1}
 \le n^{-1/20} a_{s-1}
\]
for all sufficiently large \(d\). Here \(a_s\le a_{s-1}\), \(\epsilon_n\) decays faster than any inverse power of \(n\), and \(\log_2(8/7)/2>1/20\). Take the fixed length \(s_0=400\) sweeps; iterating with sweep index jumping by two gives \(a_{s_0}\le n^{-10}\).

It follows that for any ordered list of \(k\le m\) cards from any fixed start,
\begin{equation}
 d_{\rm TV}(\text{law of their positions after }s_0\text{ sweeps},
            \text{uniform ordered distinct positions})\le\delta_n:=n^{-3} \label{special:burn-active:eq:5}
\end{equation}
for large \(d\). For each next card condition first on the \emph{paths} of the previous cards on the list, taking them as the blocked cards above. The mean total variation deviation from uniform on the remaining end slots is at most \(\sqrt{n a_{s_0}}/2\) by Cauchy-Schwarz. Conditioning instead only on the previous endpoints can only lower this averaged bound, by mixtures with the same uniform comparison given those endpoints. Sum these next-card bounds to bound the joint total variation: one can interpolate laws by taking an initial part of the list with its actual endpoint law and filling the rest uniformly without replacement. Consecutive such laws differ in total variation by at most the averaged next-card deviation (with previous endpoints actual). Since \(k\sqrt{n\cdot n^{-10}}/2\le \delta_n\), this gives \eqref{special:burn-active:eq:5}. This sequential total-variation comparison also appears in Morris's exclusion-process analysis~\cite[Lemma~12]{morris-exclusion2006}.

\end{proof}

\section{Preparation and active sweeps}\label{sec:burn}

Now designate \(m\) labels as outside cards, leaving \(m\) free cards, from any deterministic layout at a sweep boundary. We will bound the expected distance of the free-slot kernel from uniform on \(S_m\), conditional on all outside paths. Index the free slots increasingly at each boundary. By Lemma~\ref{lem:slot-kernel}, interval kernels compose by convolution, with later increments on the left, independently of the incoming free-label assignment. Each factor is computed from that interval's outside paths and starting outside layout, even when the entire outside history is revealed.

Partition the time here into chunks, each a preparation interval of \(s_0\) sweeps and then an active interval of one sweep. A subprobability is a nonnegative measure of total mass at most one. The two intervals have different tasks. After averaging over the actual preparation paths, we need joint domination by a uniform endpoint free set and an independent uniform slot permutation; this will permit a Schur conjugation average. During the active sweep, we need enough fresh coin constraints that agreement of many endpoints is unlikely. We achieve both by restricting or downweighting outcomes, obtaining conditional subprobability laws whose mean missing mass will be restored at the end.

\subsection{Truncating the preparation interval} Given the chunk's starting layout with free set \(z_0\), consider the ordered endpoint injection \(y\) at the end of the preparation interval of fictitious free labels placed canonically in the free slots at its start. Let \(p(y)\) be its probability \emph{without} conditioning on paths, from this start. Equivalently \(y\) encodes the endpoint free set \(z\) and the preparation permutation \(g\) on the canonical indices. Write \(u_*=(\binom{n}{m}m!)^{-1}\) for the uniform injection probability. Weight the preparation outcomes by \(\min(1,u_*/p(y))\) (use 1 at \(p(y)=0\)), using \(y\) of this interval's canonical increment, not the endpoints of labels carried in from prior chunks. Denote the resulting subprobability on \(g\) conditional on preparation outside paths \(e\) by \(\widetilde\nu_b^e\). Since \(z=z(e)\) is determined by those paths,
\begin{equation}
 \mathbb E_e\big[{\bf 1}_{\{z(e)=z\}}\widetilde\nu_b^e(g)\big]
 =\min(p(y),u_*)\le u_* \qquad\text{for each }(z,g). \label{special:burn-active:eq:6}
\end{equation}
Thus the retained joint measure averaged over paths is dominated by the law of uniform \(z\) and independent uniform \(g\). By \eqref{special:burn-active:eq:5} the expected removed mass is at most \(\delta_n\). Also the unweighted marginal of \(z\) is within \(\delta_n\) of uniform.

\subsection{Active truncation} For each free carrier (card considered from its slot at the start of the active sweep), count the stages of this sweep at which it is on an edge with two free slots. Call the carrier poor if this count is less than \(d/4\). Retain the active trajectory if at most
\[
 L=n^{1-1/40}
\]
free carriers are poor, dropping it otherwise. This criterion does not depend on their incoming labeling. With active outside paths \(e'\), write \(\widetilde\nu_a^{e'}\) for the conditional subprobability law of the active permutation.

At a uniform start set \(z\) independent of the fresh switches, the probability of dropping is at most \(n^{-1/40}\) for large \(d\). In fact for any full sweep switch realization form the meeting graph on the \(n\) initial slots, joining two whose carriers occupy the two slots of the same switch at some stage. It has \(d\) distinct neighbors per slot: after such a meeting at stage \(i\) the two differ in bit \(i\) and keep differing there through the sweep, preventing any later meeting in it. The graph does not use \(z\). Given a particular slot is in the uniform free set, for its \(d\) neighbors the probability any \(k\) specified ones are all outside is at most \((m/(n-1))^k\) (sampling without replacement). As in \eqref{special:burn-active:eq:3}, the probability of having more than \(3d/4\) outside neighbors is at most
\[
 (1+m/(n-1))^d 2^{-3d/4}\le n^{-1/20}
\]
for large \(d\) (the limiting \(\log_2(3/2)<7/10\)). This bounds the poor probability for that free slot. The expected number poor is thus at most \(m n^{-1/20}\); apply Markov's inequality to obtain the dropping bound. Generating full switches here computes exactly the expected removed conditional mass when averaging over the outside paths and the free switches. This path-averaged loss as a function of the active starting layout depends only on its free set (relabeling the outside has no effect). Thus from the end of the unweighted preparation interval its expectation is at most \(n^{-1/40}+\delta_n\), for any start of the chunk.

The active subprobability has zero sign transform:
\[
 \sum_g \widetilde\nu_a^{e'}(g)\operatorname{sgn}(g)=0.
\]
The free-free edge sets are determined by \(e'\). If there are none, no trajectory is retained for large \(d\) (\(m>L\)). Otherwise flip the free coin on one fixed free-free edge at the last stage having any. The counts for the carriers are unchanged (no further free-free stages), retention is invariant, and this probability-preserving involution changes the endpoint parity by exchanging the endpoints of the two carriers there.

\subsection{Trace estimate for an active interval}

Use unitary complex irreducible representations \(\rho\) of \(S_m\), writing the transform of a measure as \(F=\sum_g \nu(g)\rho(g)\). Denote by \(D\) the dimension and \(\chi\) the (unnormalized) character. We use the ordinary trace, with $\|B\|_{\HS}^2=\operatorname{tr}(B^*B)$. Write \(F_a\) for the transform of \(\widetilde\nu_a^{e'}\), and put
\[
 A(z)=\mathbb E_{e'}[F_a^*F_a\mid\text{active start with free set } z],
 \qquad b_\rho=\mathbb E_{z\ {\rm uniform}}\frac{\operatorname{tr} A(z)}{D}.
\]
The expectation defining \(A(z)\) is the same for any outside labeling of the other slots: such a relabeling just bijects outside paths with the same constraints on switch moves and the same probabilities and free kernels. For the sign representation \(F_a=0\). We treat \(D>1\) next.
The representation-theoretic background used here is recalled in~\cite{sagan,etingof}.

For the trace, conditional on \(e'\), take two independently generated free switch realizations \(X,Y\); letting \(g_X,g_Y\) be the resulting permutations, expansion gives \(\operatorname{tr}(F_a^*F_a)\) equal to the conditional average of \(\chi(g_X^{-1}g_Y)\) times the indicators that both are retained. Let \(f\) be the number of fixed points of \(g_X^{-1}g_Y\). Given a retained \(X\), for any \(a\) specified free starting indices to be fixed, their carriers must have the same paths in \(Y\) as in \(X\), by single-sweep path uniqueness from endpoints. At least \(a-L\) of them are not poor in \(X\), so they jointly visit at least \(d(a-L)/8\) distinct free-free switches (at most two carriers per switch). The coins there in \(Y\) would all have to match, and they are independent fair coins conditional on \(e'\) and \(X\). By a union bound, with \(\beta=1/100\) and for integers \(n^{1-\beta}\le a\le m\),
\begin{equation}
 \mathbb P(f\ge a,\ X,Y\text{ retained}\mid e')
 \le \binom{m}{a} 2^{-d(a-L)/8}
 \le n^{-a/32} \label{special:burn-active:eq:7}
\end{equation}
for large \(d\), since \(L\le a/2\) and \((em/a)^a\le n^{2\beta a}\) (here \(e\) in \(em/a\) is the base of natural logarithms).

The endpoint-agreement estimate must now be converted into a character estimate. We record the two bounds used here and in the reset and overlay constructions.

\begin{lemma}[Characters and fixed points]\label{lem:character-fixed-points}\label{input:character}
Fix $0<\delta<1/4$. For all sufficiently large integers $r$, every irreducible character $\chi$ of $S_r$, of degree $D=\chi(1)$, and every $\sigma\in S_r$ with $f$ fixed points satisfy
\begin{align}
 f\le r^{1-4\delta}&\quad\Longrightarrow\quad
       \frac{|\chi(\sigma)|}{D}\le D^{-\delta},\label{eq:character-small-fixed}\\
 \frac{|\chi(\sigma)|}{D}&\le D^{-\delta}r^{\delta f}.\label{eq:character-all-fixed}
\end{align}
The threshold for $r$ depends only on $\delta$.
\end{lemma}
\begin{proof}
We use the uniform Larsen--Shalev bound~\cite[Theorem~1.1]{larsen-shalev2008}, in the fixed-slack formulation of~\cite[Definition~2.1 and Theorem~2.2]{keller-lifshitz-sheinfeld2024}. If $a_i$ counts the $i$-cycles of $\sigma$, define nonnegative $e_i$ by
\[
 \sum_{i\le s}e_i=\log_r\max\left(1,\sum_{i\le s}i a_i\right),
 \qquad E(\sigma)=\sum_{i=1}^r\frac{e_i}{i}.
\]
For each fixed $\varepsilon>0$ and all sufficiently large $r$, the bound is $|\chi(\sigma)|\le D^{E(\sigma)+\varepsilon}$, uniformly in $\sigma$ and $\chi$. Since $\sum_i e_i=1$ and $e_1=\log_r\max(f,1)$,
\[
 E(\sigma)\le\frac{1+\log_r\max(f,1)}2.
\]
Taking $\varepsilon=\delta$ proves~\eqref{eq:character-small-fixed}.

For~\eqref{eq:character-all-fixed}, conjugate the fixed points to the last $f$ symbols and restrict to $S_{r-f}$. The Young branching rule gives at most $r^f$ summands counted with multiplicity, whose dimensions sum to $D$. If $r-f$ is above an absolute threshold, the remaining permutation has no fixed points, so the same character bound with $\varepsilon=1/4$ bounds its character on a summand of dimension $D_j$ by $D_j^{3/4}$. H\"older's inequality therefore gives
\[
 |\chi(\sigma)|\le \sum_j D_j^{3/4}
       \le r^{f/4}D^{3/4}.
\]
Combine the normalized estimate with $|\chi(\sigma)|/D\le1$. For every $x>0$ and $0<\delta<1/4$, $\min(1,x^{1/4})\le x^\delta$; use $x=r^f/D$. If $r-f$ is below the absolute threshold, then for large $r$ we have $f\ge r/2$ and $r^f\ge\sqrt{r!}\ge D$, so the desired bound follows from the trivial estimate. This also makes the threshold uniform.
\end{proof}

Apply the lemma with $r=m$ and $\delta=\beta/8=1/800$. For large $n$, $n^{1-\beta}\le m^{1-4\delta}$, so pairs with $f\le n^{1-\beta}$ contribute at most $D^{-\delta}$ to the normalized trace. For larger integer $f$, use~\eqref{special:burn-active:eq:7} and~\eqref{eq:character-all-fixed}; their total contribution is at most
\[
 D^{-\delta}\sum_{f>n^{1-\beta}}n^{-f/32}m^{\delta f}
 \le D^{-\delta}\sum_{f>n^{1-\beta}}n^{-(1/32-\delta)f}
 =o(1)D^{-\delta}.
\]
The estimate holds uniformly over the outside paths. In particular
\begin{equation}
 0\le b_\rho\le 2D^{-1/800}
 \qquad\text{for all sufficiently large }n.\label{special:burn-active:eq:8}
\end{equation}

To sum the trace estimates, we also need a bound on small representation degrees. The following elementary argument proves the reciprocal-square case of the more general summability theorem of Liebeck and Shalev~\cite[Theorem~2.6]{liebeck-shalev2004}. This case suffices for all the sums below.

\begin{lemma}[Reciprocal degrees]\label{lem:degree-sum}
For the irreducible degrees $D_\lambda$ of $S_r$,
\begin{equation}
 \sum_{\lambda:\,D_\lambda>1}D_\lambda^{-2}=o(1)
       \qquad(r\longrightarrow\infty).\label{eq:degree-square-sum}
\end{equation}
\end{lemma}
\begin{proof}
Recall that $D_\lambda$ counts standard Young tableaux of shape $\lambda$ and equals $r!$ divided by the product of its hook lengths~\cite{sagan,etingof}. If the first row has length $r-j$, where $1\le j\le r/4$, fill its first $j$ boxes with $1,\ldots,j$. Choose $j$ of the other $r-j$ numbers for the boxes below the first row, arrange them in one fixed standard order there, and put the remaining numbers in increasing order along the first row. Every column below that row lies among its first $j$ columns, so these fillings are standard and
\[
 D_\lambda\ge\binom{r-j}{j}.
\]
There are at most $2^j$ tail shapes, by bounding partitions by compositions. Transposition gives the same bound for a long first column. The contribution from these shapes tends to zero: $\binom{r-j}{j}\ge((r-j)/j)^j\ge3^j$ gives the summable majorant $2(2/9)^j$, while each fixed-$j$ contribution tends to zero.

For all remaining shapes, apart from the single row and single column, the maximum $L$ of a row or column length is less than $3r/4$. If $L\ge r/10$, transpose if needed and retain the first row of length $L$ and $t=\lfloor L/3\rfloor$ boxes below it, deleting corners of the tail. There are at least $t$ tail boxes because $L<3r/4$. The preceding filling argument gives at least $\binom{L}{t}$ tableaux of this subdiagram, and each extends to the original diagram by adding corners. Thus $D_\lambda$ is at least exponential in $r$. If $L<r/10$, every hook has length at most $2L<r/5$, so
\[
 D_\lambda\ge\frac{r!}{(r/5)^r},
\]
again exponential in $r$. Finally the number of partitions of $r$ is $\exp(O(\sqrt r\log r))$: specify the multiplicities of parts at most $\sqrt r$ and the list of at most $\sqrt r$ larger parts. The contribution from these remaining shapes also tends to zero.
\end{proof}

In particular,
\begin{equation}
 \sum_{\rho:\,D>1}D^{-20}=o(1),\label{special:burn-active:eq:9}
\end{equation}
and the minimum nontrivial degree $D>1$ tends to infinity. Absorbing the factor $2$ in~\eqref{special:burn-active:eq:8}, we may therefore use $b_\rho\le D^{-1/1600}$ for all sufficiently large $m$.

\subsection{From traces to contraction over chunks}

The active interval has supplied an averaged normalized trace. Its matrix need not be central, and matrices from consecutive intervals need not commute. We now use preparation domination to upgrade the trace bound to a positive-operator inequality, uniformly in the incoming outside layout. This is the step that makes the conditional law of the entire assignment accessible.

Let \(F_b\) be the transform of the preparation subprobability conditional on paths \(e\). Conditional on the chunk paths the chunk subprobability (composing preparation then active) has transform \(F_a F_b\). Averaging over chunk paths from any start, its squared-matrix estimate in positive semidefinite order is
\begin{equation}
\begin{aligned}
 \mathbb E[(F_a F_b)^*(F_a F_b)]
 &=\mathbb E_e[F_b^* A(z(e)) F_b]\\
 &\preceq \mathbb E_e\sum_g\widetilde\nu_b^e(g)\rho(g)^* A(z(e))\rho(g)
 \preceq b_\rho I.
\end{aligned}\label{special:burn-active:eq:10}
\end{equation}
For the equality, the next outside paths start fresh and \(A(z)\) does not require their particular starting labeling. The first inequality is weighted Cauchy-Schwarz applied to \(A^{1/2}F_b\) on each vector, with weights of sum at most 1. For the second, each conjugated matrix is positive semidefinite, so use \eqref{special:burn-active:eq:6} to bound by the average with uniform \(z,g\). For each \(z\), the uniform \(g\) conjugation gives \(\operatorname{tr}(A(z))I/D\) by Schur's lemma.

Use \(J=40000\) chunks. Conditional on their outside paths let \(\nu\) be the full probability law of the overall free-slot permutation and let \(\widetilde\nu\) be the composition of the subprobabilities above. Then \(\widetilde\nu\le\nu\) coefficientwise. Indeed all truncations use only the randomness of their respective interval conditional on paths and the layout information there, computing increments on canonical slots regardless of previous free permutations. Write \(F_{\rm tot}\) for the transform of \(\widetilde\nu\), the product of the chunk transforms (latest on the left). Applying \eqref{special:burn-active:eq:10} given previous paths for the last chunk and iterating gives
\[
 \mathbb E[F_{\rm tot}^*F_{\rm tot}]\preceq b_\rho^J I
\]
from any starting layout: the preceding product is determined by the preceding outside paths and the bound \eqref{special:burn-active:eq:10} is uniform conditional on them. The sign transform is zero. The mass \(M\) of \(\widetilde\nu\) is the product of the interval masses. By bounding the deficit of the product by the sum of deficits,
\[
 \mathbb E[1-M]\le J(2\delta_n+n^{-1/40}).
\]
Here we average the deficits under the ordinary outside path law; the preparation and active loss bounds hold given the outside layout starting each chunk, with the unweighted preparation path law for bounding the active loss.

Finite-group Plancherel and Cauchy-Schwarz now give, with \(U_m\) uniform on \(S_m\),
\[
 \frac12\mathbb E\|\widetilde\nu-MU_m\|_1
 \le \frac12\left(\sum_{\rho:\,D>1}D^2 b_\rho^J\right)^{1/2}=o(1).
\]
Indeed for the difference \(v\) we have \(\|v\|_1^2\le |S_m|\sum_g |v(g)|^2=\sum_\rho D\|\sum_g v(g)\rho(g)\|_{\rm HS}^2\) (the equality also follows directly by the regular character identity). The trivial transform of the difference vanishes, and at nontrivial representations the uniform transform vanishes. Take expectations using the matrix bound (and bound the mean norm by its root mean square), then use \eqref{special:burn-active:eq:9} since \(J/1600=25\). Restoring the missing nonnegative measure, we conclude
\begin{equation}
 \mathbb E\ d_{\rm TV}(\nu,U_m)=o(1) \label{special:burn-active:eq:11}
\end{equation}
uniformly in the outside starting layout: add at most \(\mathbb E[1-M]\) to the centered subprobability bound by comparing \(\nu-\widetilde\nu\) to \((1-M)U_m\). Here \(\nu\) itself is conditional on all the outside paths.

\subsection{Completing the first route}\label{sec:burn-completion}

Couple a chain from any fixed deck and one started uniform (uniformity is preserved by the independent random positional permutations). After \(s_0\) sweeps we can couple the ordered positions of the \(m\) designated outside labels to agree with probability at least \(1-\delta_n\), by \eqref{special:burn-active:eq:5} and maximal coupling; complete the two endpoint decks by their respective conditional laws. On agreement sample common outside paths for the next \(J\) chunks from their path law. Conditional on these paths the final free orders in each deck have distributions obtained using \(\nu\) on their possibly different incoming free orders. Each is within \(d_{\rm TV}(\nu,U_m)\) of uniform on the same set of end free orders (permuting from any fixed incoming order). Maximally couple them, so the expected failure probability on this continuation is at most \(2\mathbb E\,d_{\rm TV}(\nu,U_m)=o(1)\) by \eqref{special:burn-active:eq:11}, uniformly on agreement layouts. This couples endpoints with the correct marginals by path conditioning (on initial disagreement one can just continue independently). The total failure probability is at most \(\delta_n+o(1)<1/4\) for sufficiently large \(d\), bounding the worst-case total variation.

The conditional assertion of Theorem~\ref{thm:methods} was established in~\eqref{special:burn-active:eq:11}, over $40000(400+1)=16040000$ sweeps. The coupling proves its full-deck conclusion: including the initial preparation gives $400+40000(400+1)=16040400$ sweeps, each costing $d$ physical shuffles. The conditional expectation is over the ordinary law of the outside paths; neither the argument nor its conclusion asserts a bound for each fixed environment.
